%=====================================================================
% Cerceveden bagimsiz ANA taslak (Ingilizce).
% ELECO surumu (Turkce + Ingilizce) papers/eleco/ altinda bundan turetilecek.
% Derleme: ./papers/build.sh draft-en   (tectonic) veya Overleaf.
%=====================================================================
\documentclass[conference]{IEEEtran}
\usepackage{cite}
\usepackage{amsmath,amssymb}
\usepackage{graphicx}
\usepackage{booktabs}
\usepackage{array}
\usepackage{url}
\graphicspath{{figures/}}

\begin{document}

\title{A Physically Faithful Device Simulator for a\\
Third-Party GNSS-Disciplined Rubidium Oscillator:\\
Construction and Fidelity Validation from Limited Data}

\author{\IEEEauthorblockN{Omer Faruk Karagoz}
\IEEEauthorblockA{Kurum adi\\ ofkaragoz05@gmail.com}}

\maketitle

%---------------------------------------------------------------------
\begin{abstract}
Software that controls or monitors a precision timing instrument must be
developed against the instrument's command interface, yet the physical
device is typically expensive, singular, slow to reach a usable state,
and unavailable for the extended periods that testing demands. A
software stand-in that speaks the same interface removes this
dependency, but raises two questions that are sharpened when the device
is a third-party product whose firmware is inaccessible: how faithful
must the stand-in be, and how can its fidelity be established at all. We
report the construction and validation of such a stand-in for a
commercial Low Noise Rubidium GNSS-disciplined oscillator. In the
software-testing taxonomy the artefact is a \emph{fake}, and in the
digital-twin taxonomy a \emph{digital model}; it is explicitly not a
digital twin, having no automated link to any physical unit. Our
contribution is not the underlying clock model, which is standard, but
a repeatable method: every observable is traced to the vendor manual or
to a small set of protocol recordings, every unsourced value is marked
as an assumption, and fidelity is argued through physical invariants,
through the reproduction of behaviours that were never tuned for, and
through documented manual-versus-device discrepancies. The resulting
model reproduces the manual's own worked health-status example without
having been fitted to it.
\end{abstract}

\begin{IEEEkeywords}
device simulator, test double, digital model, GNSS-disciplined
oscillator, holdover, SCPI, software testing, reproducibility.
\end{IEEEkeywords}

%=====================================================================
\section{Introduction}
Global Navigation Satellite System (GNSS) disciplined oscillators
(GNSSDOs) deliver traceable time and frequency to telecommunications,
power, instrumentation, and defence systems. A GNSSDO combines the
long-term accuracy of GNSS with the short-term stability of a local
oscillator; in the device studied here that oscillator is a rubidium
(Rb) atomic reference, steered by a software control loop to a
GNSS-derived reference. Any application built on such a device---a
monitoring dashboard, an automated test system, a logging or alarming
service---must be implemented and tested against the device's serial
command interface.

This dependency is awkward in practice. The instrument is costly and
often singular; it must warm up and then discipline its oscillator
before it reports usable values, a process of tens of minutes; and it
cannot be monopolised by a software team for the days of repeated,
fault-injecting runs that thorough testing requires. Worse, several of
the behaviours a client most needs to handle---loss of GNSS, holdover,
health degradation---are precisely the ones that are hardest to trigger
on demand against real hardware.

The established response is to substitute, during development, a
software stand-in that speaks the same interface, so that unmodified
client code cannot tell the difference. This turns the problem into two
engineering questions. First, \emph{how faithful} must the stand-in be:
which behaviours must it reproduce, and to what precision, for it to be
useful for its intended purpose? Second, and more delicate when the
device is a third-party product whose firmware cannot be inspected and
whose single unit is only intermittently available, \emph{how can that
fidelity be established} without the exhaustive hardware comparison that
the literature usually assumes?

This paper answers both questions for a commercial Jackson Labs Low
Noise Rubidium GNSDO. Our contributions are:
\begin{itemize}
\item \textbf{A layered, physically grounded simulator}
(Section~\ref{sec:arch}--\ref{sec:models}) that answers the device's
serial/SCPI interface and reproduces its time-dependent
behaviour---holdover error accumulation, deterministic measurement
noise, and warm-up---and that runs over a virtual serial port so that
unmodified client software connects to it as to the real device.
\item \textbf{A construction method for a third-party device}
(Section~\ref{sec:models}): every observable value is derived either
from the vendor manual or from a small set of protocol recordings, and
every value that has no such source is explicitly labelled an
assumption and exposed as configuration.
\item \textbf{A fidelity-validation method for limited access}
(Section~\ref{sec:validation}): in the absence of exhaustive hardware
comparison, fidelity is argued through automated physical-invariant
tests, through the reproduction of behaviours that were not tuned for,
and through a catalogue of manual-versus-device discrepancies that a
documentation-only simulator would get wrong.
\end{itemize}

We are deliberately precise about what the artefact is. By the
digital-twin taxonomy it is a \emph{digital model}, and in
software-testing terms a \emph{fake}; it is \emph{not} a digital twin.
Section~\ref{sec:related} justifies this terminology, which we regard as
a matter of scientific accuracy rather than modesty.

%=====================================================================
\section{Background}
\label{sec:background}
\subsection{Device and interface}
The reference device is a Jackson Labs Low Noise Rubidium
GNSDO~\cite{jacksonlabs2015lnrb}. It exposes a serial interface at
$115{,}200$~baud, 8N1, using the SCPI convention: ASCII commands
terminated by \texttt{CR}, \texttt{LF}, or \texttt{CR}/\texttt{LF},
case-insensitive, with both abbreviated and long spellings
(e.g.\ \texttt{SYNC:TINT?} and \texttt{SYNChronization:TINTerval?}).
Queries end in \texttt{?}; a smaller set of commands change device
state.

\subsection{Disciplining and the time-interval error}
Two one-pulse-per-second (1PPS) signals coexist inside the unit. One is
produced by the GNSS receiver: accurate over the long term because it
is tied to GNSS system time, but noisy from second to second. The other
is produced by the Rb oscillator: extremely quiet over the short term,
but subject to slow drift. A control loop steers the Rb oscillator so
that its 1PPS tracks the GNSS 1PPS. The instantaneous difference between
the two, the \emph{time-interval error} (TINT, queryable as
\texttt{SYNC:TINT?}), is the loop's error signal and is reported with a
resolution of $10^{-10}$~s.

\subsection{Holdover}
While the device is locked, the closed loop drives TINT towards zero. If
GNSS is lost the loop opens: no correction is applied, and the
oscillator's residual frequency error integrates into a growing
\emph{time} error. This coasting regime is called \emph{holdover}. Its
error growth---and the resulting degradation of the device's reported
health---is the first and most safety-relevant behaviour a client must
handle, and the first we model.

\subsection{Health status and warm-up}
The device summarises its condition in a hexadecimal bit-mask
(\texttt{SYNC:HEALTH?}), each bit encoding one fault or transient
condition (for example, phase offset beyond $210$~ns, or holdover
longer than $60$~s). From a cold start the unit is not immediately
usable: the atomic reference reaches operating temperature in under two
minutes, but full GNSS lock takes on the order of twenty minutes, and
the device advances through distinct servo states in between. These
timescales matter to client software and are modelled in
Section~\ref{sec:models}.

%=====================================================================
\section{Related Work and Positioning}
\label{sec:related}
\subsection{Digital models, shadows, and twins}
The term ``digital twin'' is applied so broadly that it has become
ambiguous. We adopt the categorical distinction of
Kritzinger~\emph{et~al.}~\cite{kritzinger2018dt}, formalised by Barbie
and Hasselbring~\cite{barbie2024dtp}, which classifies a virtual
counterpart by the degree of \emph{automated} data exchange with a
specific physical instance: a \emph{digital model} has no automated
link; a \emph{digital shadow} has an automated one-way
physical-to-digital link; a \emph{digital twin} has an automated,
bidirectional link over the instance's life cycle. Our artefact has no
automated link of any kind---its parameters were extracted, once and by
hand, from a manual and a handful of recordings---so it is, precisely, a
digital model. The interactive command/response exchange between a test
client and the simulator is unrelated to this classification, which
concerns the link between a physical instance and its digital
counterpart.

Within the same body of work, Barbie and Hasselbring define a
\emph{Digital Twin Prototype} (DTP): a software prototype whose emulated
sensors are driven by prior recordings, used for continuous-integration
testing without the physical
object~\cite{barbie2024dtp,barbie2023smartfarming}. Our purpose and
architecture are close to the DTP idea, down to keeping the serial
protocol at the host layer so that device drivers are unchanged. The
decisive difference is one of access and provenance: a DTP presumes the
device's \emph{own} firmware (``the configurations are equal''), whereas
we reconstruct behaviour from outside the protocol boundary, black-box,
for a device we did not design and cannot open. Khedr~\emph{et~al.}
observe that digital-twin engineering still lacks practical guidance for
formalising and verifying qualities such as \emph{substitutability} and
\emph{fidelity}~\cite{khedr2026sos}; these are exactly the qualities our
validation method addresses, in the specific and common setting of a
third-party instrument.

\subsection{Test doubles}
In the vocabulary of software testing~\cite{meszaros2007xunit,%
fowler2007mocks} the artefact is a \emph{fake}: a working but simplified
implementation of a depended-on component, as distinct from a
\emph{stub} or \emph{mock}, which return pre-arranged answers without
modelling behaviour. This distinction also describes the project's own
history: an earlier version returned constant values and was, in effect,
a stub; the work reported here made it a fake. Deriving a double's
behaviour from recordings of the real system is the principle behind
RICK~\cite{tiwari2023rick}, which generates mocks from production
traces; we apply the same principle manually, to device protocol
recordings.

\subsection{Oscillator and holdover modelling}
Modelling oscillator stability and holdover is a mature area, with
refinements using Kalman filtering, machine learning, and explicit
quantisation-noise models. We make \emph{no} claim of contribution
there: we deliberately use the standard two-state clock model with
parameters taken from the vendor manual. We also distinguish our work
from GNSS \emph{signal} simulators, which synthesise radio-frequency
satellite signals; our artefact is a device- and protocol-level
simulator that synthesises SCPI responses. The contribution lies in the
method by which the model is built and validated, not in the model
itself.

%=====================================================================
\section{Simulator Architecture}
\label{sec:arch}
The simulator is organised in layers so that device physics is isolated
from protocol handling (Fig.~\ref{fig:arch}):
\begin{itemize}
\item a \emph{serial server} that owns a pseudo-terminal, created with
\texttt{socat}, and converts between the byte stream and lines; a
client opens this virtual port exactly as it would a real
\texttt{/dev/tty} device;
\item an \emph{SCPI parser} that normalises a command and dispatches it
to a registered handler, resolving abbreviated and long forms, and
returning \texttt{Command Error} for unrecognised input, as the real
device does;
\item a \emph{device-state} object that holds the instrument's memory
and is the single point from which all responses are derived, so that a
summary query and an individual query can never disagree;
\item a \texttt{models} package of \emph{pure functions} encoding the
physics, performing no input/output and receiving elapsed time as an
argument.
\end{itemize}
The purity of the model layer is a deliberate design choice with a
direct bearing on validation (Section~\ref{sec:validation}): because a
model function is a function of elapsed time alone, the question ``what
would the device read after five hours of holdover?'' is answered in a
unit test in microseconds, without waiting and without a running server.
New command groups are added by registering handlers, so the protocol
surface grows without changes to the core.

\begin{figure}[t]
\centering
\includegraphics[width=\columnwidth]{fig13_architecture.pdf}
\caption{Layered architecture. Unmodified client software connects
through a virtual serial port; responses are derived from a single
device-state object, which in turn reads the pure-function physical
models.}
\label{fig:arch}
\end{figure}

%=====================================================================
\section{Physical Models}
\label{sec:models}
Every numeric parameter below is traced to a source in
Table~\ref{tab:params}; the single parameter with no documentary source,
the ambient start temperature, is marked as an assumption and is
configurable.

\subsection{Holdover error accumulation}
Phase is the integral of frequency. Let $y_0$ be the fractional
frequency error at the instant GNSS is lost and let $D$ be the
oscillator's ageing rate. The accumulated time error is
\begin{equation}
x(t) = x_0 + y_0\,t + \tfrac{1}{2}\,D\,t^{2},
\label{eq:holdover}
\end{equation}
where $x_0$ is the phase frozen at the moment of loss, since phase is
continuous and does not jump to zero. The linear term dominates over
minutes-to-hours; the quadratic ageing term is negligible there and
becomes visible only over days (Fig.~\ref{fig:holdover}). We take $y_0$
to be the device's own Frequency Error Estimate
(\texttt{SYNC:FEE?})---the manual defines it as an Allan-variance-like
figure over a $1000$~s interval---and $D$ to be the manual's
$8\times10^{-14}$/day ageing specification. Because $y_0$ is read from
the device, the accumulation rate is consistent with what the device
reports, and a client can in principle predict its own holdover budget
(Fig.~\ref{fig:fee}).

\subsection{Deterministic measurement noise}
A real instrument's temperature, voltage, and TINT readings are never
perfectly constant, but modelling their variation as an independent
random draw per query would be wrong on two counts. Physically, the
device has a definite value at a given instant: two queries at the same
time must return the same number, whereas an independent draw would make
the value depend on how often it is polled. Methodologically, an
independent draw is not reproducible, which frustrates debugging. We
therefore model each reading as a deterministic function of elapsed
time,
\begin{equation}
v(t) = v_{\mathrm{base}} + A\, f(t),\qquad
f(t)=\!\!\sum_{k}\! w_k \sin\!\big(2\pi \phi^{k} t/P + \varphi_k\big),
\label{eq:noise}
\end{equation}
where the frequency ratios are powers of the golden ratio $\phi$, so
that the composite pattern is quasi-periodic and does not visibly
repeat. This yields both required properties: identical values for
identical instants, and exact run-to-run repeatability
(Fig.~\ref{fig:determinism}). Amplitudes $A$ are taken from observed
ranges and the periods $P$ are separated by physical timescale---slow
for temperature, which has thermal inertia, faster for electrical
quantities.

\subsection{Warm-up and servo states}
From a cold start the temperature approaches its steady value
exponentially,
\begin{equation}
T(t) = T_\infty - (T_\infty - T_0)\,e^{-t/\tau},
\label{eq:warmup}
\end{equation}
because heat loss scales with the temperature difference; a linear ramp
would be unphysical, overshooting rather than asymptoting. The thermal
rise is fast---of the order of a few minutes, consistent with the
manual's ``less than two minutes to atomic lock''---whereas GNSS
\emph{lock} is a separate and slower process of roughly twenty minutes.
The device therefore advances through servo states
(\texttt{SERV:STATE?}) from warm-up, to locking, to locked
(Fig.~\ref{fig:warmup}); \texttt{SYNC:LOCKED?} alone reports only the
last of these transitions and cannot distinguish ``still warming'' from
``warm but not yet locked''.

\begin{table}[t]
\caption{Traceability of model parameters}
\label{tab:params}
\centering
\renewcommand{\arraystretch}{1.2}
\begin{tabular}{@{}p{0.30\columnwidth}p{0.30\columnwidth}p{0.30\columnwidth}@{}}
\toprule
\textbf{Parameter} & \textbf{Value} & \textbf{Source} \\
\midrule
Ageing $D$ & $8\times10^{-14}$/day & Manual, \S2.9 \\
$y_0$ (FEE) & device-reported & Manual, \S3.6.10; recording \\
Locked TINT & $\approx 0$, $\pm12$~ns jitter & Recordings \\
CSAC$-$PCB offset & $+1.32\,^\circ$C & Recordings (4 pairs) \\
EFC scale & $\times 200$ (ppt/\%) & Recordings; Manual \S3.7 \\
Phase threshold & $210$~ns & Manual, \S3.6.18 \\
Thermal $\tau$ & $\approx 90$~s & Operator observation; \S1.1 \\
Ambient $T_0$ & $25\,^\circ$C & \emph{Assumption (configurable)} \\
\bottomrule
\end{tabular}
\end{table}

%=====================================================================
\section{Fidelity Validation}
\label{sec:validation}
The central methodological question of this work is how to argue that a
simulator is faithful when the device it imitates is a third-party
product available only intermittently, so that the exhaustive
side-by-side comparison assumed by much of the literature is
impossible. We use four complementary strategies.

\subsection{Traceability of parameters}
Every parameter is tied to a documentary source---a manual clause or a
protocol recording (Table~\ref{tab:params}). The one parameter without
such a source is labelled an assumption and exposed as configuration, so
that a disputed value is a single, visible edit rather than a constant
buried in code. This makes the model's empirical commitments auditable.

\subsection{Physical-invariant testing}
Each behaviour is exercised by an automated measurement that asserts the
physical invariants it must satisfy---holdover monotonic and signed
correctly, locked TINT zero-mean and trend-free, temperature rising and
bounded, repeated runs identical---and that writes its data to a
versioned artefact stamped with the source-control revision. Because the
model functions are pure, these measurements sweep hours or days of
simulated time instantly. The figures in this paper are regenerated from
those artefacts by a single command, so every reported result is
reproducible from the recorded data.

\subsection{Reproduction of untuned behaviour}
The strongest available evidence that a model is correct, short of full
hardware comparison, is that it reproduces an observation for which it
was never tuned. Three such coincidences arose, none of them targets of
the fitting process. First, after a multi-hour holdover the simulator
returns the health status \texttt{0x14}; this is exactly the worked
example given in the manual (\texttt{0x10}\,$\vert$\,\texttt{0x4}, i.e.\
``in holdover'' combined with ``phase beyond $210$~ns''),
Fig.~\ref{fig:health}. Second, the per-second growth rate of TINT in
holdover equals the device's independently reported FEE. Third, the CSAC
temperature tracks the PCB temperature at a constant offset that matches
all four recorded pairs to within $0.07\,^\circ$C
(Fig.~\ref{fig:csac}). We present these as corroborating evidence, not
proof; but a model built by fitting other quantities has no reason to
reproduce them unless its underlying structure is right.

\subsection{Documented manual-versus-device discrepancies}
Building from recordings, rather than from the manual alone, exposed
several points where the documentation is insufficient or misleading
(Table~\ref{tab:discrepancies}). Each is a case in which a
documentation-only simulator would be wrong, and each is therefore
evidence that faithful behaviour must be aligned to observation.

\begin{table}[t]
\caption{Manual-versus-device discrepancies found during construction}
\label{tab:discrepancies}
\centering
\renewcommand{\arraystretch}{1.25}
\begin{tabular}{@{}p{0.30\columnwidth}p{0.60\columnwidth}@{}}
\toprule
\textbf{Observable} & \textbf{Discrepancy} \\
\midrule
\texttt{SYNC:TINT?} & Signed, swinging about zero; the manual's
``0.2\,ns average'' is a mean, not an amplitude. \\
\texttt{MEAS:CURR?} & Returns a temperature, not a current (legacy
command). \\
\texttt{EFControl} & Absolute equals relative $\times 200$; relation not
documented. \\
Satellite counters & Tracked count can exceed visible count in
recordings. \\
\bottomrule
\end{tabular}
\end{table}

%=====================================================================
\section{Results}
\label{sec:results}
All figures are produced by the two-stage measurement/plot pipeline of
Section~\ref{sec:validation}: a measurement stage runs the model,
asserts invariants, and logs a data artefact; a plotting stage reads
that artefact. Fig.~\ref{fig:holdover} shows holdover accumulation from
\eqref{eq:holdover}: the error crosses the $210$~ns health threshold at
about $3.25$~hours for the device's reported FEE, and the quadratic
ageing contribution (shaded) is negligible over this span.
Fig.~\ref{fig:beforeafter} contrasts this with the project's earlier
static behaviour, a constant $850$~ns, whose shaded difference is the
error that a client would have been told to ignore.
Fig.~\ref{fig:health} traces the health register through holdover,
reaching the manual's \texttt{0x14}. Fig.~\ref{fig:warmup} shows the
fast thermal rise against the slow arrival of GNSS lock, and
Fig.~\ref{fig:csac} shows the temperature-offset fidelity check against
the recorded pairs. Further figures produced by the same
pipeline---deterministic measurement noise, run-to-run determinism, the
EFC relation, the divergence of the main and filter time-interval loops
in holdover, FEE sensitivity, and the emulated device's Allan-deviation
signature---are available and selected per venue.

%=====================================================================
\section{Limitations}
The artefact is a digital model, not a shadow or a twin: it has no live
link to any physical unit, and does not update as the device ages. The
measurement noise is deterministic and does \emph{not} reproduce the
device's spectral, Allan-deviation signature; this is a deliberate
trade-off, since the intended purpose---testing client software---values
repeatability over spectral realism, but it means the simulator must not
be used to evaluate oscillator stability. Some recorded values were
untimestamped, so certain behaviours are validated as ranges and
relationships rather than as time-anchored points. Finally, validation
rests on a single device instance and a limited set of recordings;
broader hardware comparison remains future work.

%=====================================================================
\section{Conclusion and Future Work}
We have built a physically grounded, behaviourally validated simulator
of a third-party GNSS-disciplined rubidium oscillator from only its
vendor manual and a small set of protocol recordings, and we have set
out a method for arguing fidelity when access to the device is limited:
trace every parameter, test physical invariants automatically, seek the
reproduction of untuned behaviour, and document where the manual and the
device disagree. A natural next step, and the one that would change the
artefact's taxonomic status, is to let the simulator ingest periodic
device logs and self-calibrate, establishing the automated one-way link
that would make it a digital \emph{shadow}.

%---------------------------------------------------------------------
\bibliographystyle{IEEEtran}
\bibliography{references}

\end{document}
